% !TeX root = ../main.tex

% ==================================================
% 2.1 Introduction
% ==================================================

\subsection{Introduction - 클래스와 객체 모델을 왜 공부하는가?}

\begin{dhknote}
\begin{dhksmry}[centered]
학습목표
\end{dhksmry}
\end{dhknote}

\begin{dhksmry}[centered]
Python의 클래스와 객체 모델을 공부해야 하는 이유를

이해하고 핵심 질문을 정의한다.
\end{dhksmry}

이번 학습의 출발점은 다음과 같은 PyTorch 모델 코드다.

\begin{dhkcode}
import torch.nn as nn


class MyModel(nn.Module):

    def __init__(self):
        super().__init__()
        self.linear = nn.Linear(10, 5)

    def forward(self, x):
        return self.linear(x)
\end{dhkcode}

이 코드는 PyTorch를 이용해 간단한 신경망 모델을 정의한 것이다.

코드의 기능을 기준으로 설명하면 다음과 같다.

\begin{dhksmry}
    1. \zsh{MyModel}이라는 클래스를 정의한다.

    2. \zsh{nn.Module}을 상속한다.

    3. \zsh{__init__()}에서 신경망 계층을 정의한다.

    4. \zsh{forward()}에서 입력 데이터에 대한 연산을 정의한다.
\end{dhksmry}

이 정도의 설명만으로도 코드의 목적은 대략적으로 파악할 수 있다.

하지만 코드의 기능을 이해하는 것과 Python이 해당 코드를 실행하는 원리를 이해하는 것은 서로 다른 문제다.

예를 들어 다음과 같은 질문을 생각해 볼 수 있다.

\begin{dhktable}
|코드	|핵심 질문|
|:---|:---|
|class MyModel:	|클래스 정의를 실행하면 무엇이 생성되는가?|
|MyModel()	|클래스 이름 뒤에 ()를 붙이면 어떤 과정이 시작되는가?|
|__init__(self)	|왜 인스턴스를 생성할 때 자동으로 호출되는가?|
|self	|무엇을 참조하며 누가 전달하는가?|
|self.linear = ...	|생성한 객체는 어디에 저장되는가?|
|nn.Module	|상속하면 어떤 기능을 사용할 수 있는가?|
|super().__init__()	|상위 클래스의 초기화가 왜 필요한가?|
|model(x)	|모델 인스턴스를 함수처럼 호출할 수 있는 이유는 무엇인가?|
|__call__()	|Python은 왜 이 이름을 특별하게 처리하는가?|
\end{dhktable}

이러한 질문에 순서대로 답하는 것을 목표로 한다.

\subsubsection{메서드와 \zsh{self}에 대한 의문}

다음 코드를 살펴보자.

\begin{dhkcode}
class Person:

    def introduce(self):
        print("Hello")


p = Person()

p.introduce()
\end{dhkcode}

실행 결과:

\begin{dhkzsh}
Hello
\end{dhkzsh}

여기서 \zsh{introduce()} 메서드는 다음과 같이 정의되어 있다.

\begin{dhkcode}
def introduce(self):
\end{dhkcode}

분명히 \zsh{self}라는 매개변수가 하나 존재한다.

그런데 실제 호출에서는 아무런 인수도 전달하지 않았다.

\begin{dhkcode}
p.introduce()
\end{dhkcode}

그럼에도 정상적으로 실행된다.

왜 그럴까?

Python에서는 인스턴스를 통해 일반적인 인스턴스 메서드에 접근하면 해당 인스턴스가 연결된 \zsh{바운드 메서드(Bound Method)}를 얻기 때문이다.

따라서 다음 두 호출은 이 예제에서 같은 효과를 가진다.

\begin{dhkcode}
p.introduce()

Person.introduce(p)
\end{dhkcode}

여기서 \zsh{self}는 Python의 예약어가 아니다.

메서드에 전달되는 인스턴스를 받기 위해 관례적으로 사용하는 매개변수 이름이다.

즉, \zsh{self}를 이해하려면 다음 개념을 함께 이해해야 한다.

\begin{dhksmry}
    1. 함수도 객체라는 사실

    2. 클래스에 정의된 함수와 인스턴스 메서드의 관계

    3. 인스턴스를 통한 속성 접근

    4. 바운드 메서드와 인수 전달
\end{dhksmry}

이 내용은 \zsh{2.5. 메서드와 self}에서 자세히 학습한다.

\subsubsection{인스턴스 생성과 초기화에 대한 의문}

다음 코드를 살펴보자.

\begin{dhkcode}
class Person:
    def __init__(self, name):
        self.name = name

p = Person("Kim")
\end{dhkcode}

일반적으로 \zsh{__init__()}은 객체를 생성할 때 실행되는 함수라고 설명한다.

하지만 엄밀하게 말하면 \zsh{__init__()}은 객체 자체를 생성하는 메서드가 아니다.

Python의 일반적인 인스턴스 생성 과정은 다음과 같다.

\begin{dhksmry}[centered]
Person("Kim")
\da
클래스 객체 호출
\da
__new__()
\da
인스턴스 생성
\da
__init__()
\da
인스턴스 초기화
\da
p에 인스턴스 바인딩
\end{dhksmry}

이것은 일반적인 클래스 생성 과정의 개념도이며, 

실제로는 \zsh{__new__()}가 해당 클래스의 인스턴스를 반환하는지 등에 따라

\zsh{__init__()} 호출 여부가 달라질 수 있다.

여기서 두 메서드의 역할을 구분해야 한다.

\begin{dhktable}
|메서드	|주요 역할|
|:---|:---|
|__new__()	|인스턴스 생성|
|__init__()|	생성된 인스턴스 초기화|
\end{dhktable}

그리고 다음 코드에서는

\begin{dhkcode}
self.name = name
\end{dhkcode}

인스턴스에 \zsh{name}이라는 속성을 설정한다.

따라서 다음 질문들이 발생한다.

\begin{dhksmry}
    1. __new__()와 __init__()은 왜 분리되어 있는가?

    2. Python은 이 메서드들을 어떤 규칙으로 호출하는가?

    3. self는 어느 시점에 전달되는가?

    4. self.name은 어디에 저장되는가?

    5. __init__()은 어떤 값을 반환해야 하는가?
\end{dhksmry}

이 내용은 \zsh{2.6. __new__와 __init__}에서 자세히 학습한다.

\subsubsection{특수 메서드와 Callable 객체}

다음 코드를 살펴보자.

\begin{dhkcode}
class Calculator:
    def __call__(self, x):
        return x * 2

calc = Calculator()

result = calc(10)

print(result)
\end{dhkcode}

실행 결과:

\begin{dhkzsh}
20
\end{dhkzsh}

여기서 \zsh{calc}는 함수가 아니다.

\zsh{Calculator} 클래스의 인스턴스다.

그런데 다음 코드가 정상적으로 실행된다.

\begin{dhkcode}
calc(10)
\end{dhkcode}

이것이 가능한 이유는 \zsh{Calculator} 클래스에 \zsh{__call__()}이라는 특수 메서드가 정의되어 있기 때문이다.

Python은 인스턴스가 호출될 때 해당 타입의 \zsh{__call__()}을 이용한다.

따라서 이 예제에서는 다음과 같은 호출 관계를 생각할 수 있다.

\begin{dhksmry}[centered]
calc(10)
\da
Calculator.__call__(calc, 10)
\da
10 * 2
\da
20
\end{dhksmry}

이러한 객체를 Callable 객체라고 한다.

이 개념은 PyTorch 모델에서도 사용된다.

\begin{dhkcode}
model = MyModel()

output = model(x)
\end{dhkcode}

여기서 \zsh{model}은 \zsh{MyModel}의 인스턴스다.

하지만 \zsh{nn.Module}이 제공하는 호출 메커니즘을 상속받았기 때문에 함수처럼 호출할 수 있다.

PyTorch는 이 호출 과정에서 \zsh{forward()}를 실행하도록 구성되어 있다.

따라서 \zsh{__call__()}을 이해하면 다음 두 코드의 관계도 설명할 수 있다.

\begin{dhkcode}
model(x)

model.forward(x)
\end{dhkcode}

두 코드 모두 순전파 계산을 수행할 수 있지만, 동일한 호출 경로를 거치는 것은 아니다.

\zsh{model(x)}는 PyTorch의 모듈 호출 메커니즘을 거치므로 일반적으로 권장되는 방식이다.

이 내용은 \zsh{2.10. 특수 메서드}에서 자세히 학습한다.

\subsubsection{상속은 왜 사용하는가?}

다음 코드를 살펴보자.

\begin{dhkcode}
class MyModel(nn.Module):
    ...
\end{dhkcode}

여기서 \zsh{MyModel}은 \zsh{nn.Module}을 상속한다.

상속은 일반적으로 기존 클래스의 기능을 재사용하는 방법으로 소개된다.

하지만 상속의 목적은 단순한 코드 재사용에만 있지 않다.

상속을 통해 기존 클래스가 제공하는 메서드와 동작 규약을 활용하거나,

필요한 부분을 재정의하여 새로운 클래스를 구성할 수 있다.

예를 들어 \zsh{nn.Module}을 상속받은 모델에서는 다음과 같은 메서드들을 사용할 수 있다.

\begin{dhkcode}
model.parameters()

model.state_dict()

model.to(device)

model.train()

model.eval()
\end{dhkcode}

우리는 이러한 메서드를 \zsh{MyModel} 내부에 직접 정의하지 않았다.

그런데도 사용할 수 있는 이유는 상속 관계를 통해 해당 기능들이 제공되기 때문이다.

또한 \zsh{nn.Module}은 하위 모듈과 파라미터를 등록하고 관리하는 기능도 제공한다.

따라서 다음 코드에서

\begin{dhkcode}
self.linear = nn.Linear(10, 5)
\end{dhkcode}

\zsh{nn.Linear} 인스턴스를 모델의 하위 모듈로 등록하는 동작이 이루어진다.

이러한 구조를 이해하려면 다음 개념이 필요하다.

\begin{dhksmry}
    1. 상속 관계

    2. 메서드 탐색

    3. 메서드 오버라이딩

    4. super()

    5. MRO(Method Resolution Order)

    6. 클래스가 제공하는 객체 관리 기능
\end{dhksmry}

이 내용은 \zsh{2.8. 상속과 메서드 오버라이딩}, \zsh{2.9. super()와 MRO}에서 자세히 학습한다.

\subsubsection{최종 학습 목표}

이번 학습의 최종 목표는 Python 클래스의 문법을 단순히 암기하는 것이 아니다.

다음 코드를 다시 살펴보자.

\begin{dhkcode}
import torch.nn as nn


class MyModel(nn.Module):
    def __init__(self):
        super().__init__()
        self.linear = nn.Linear(10, 5)

    def forward(self, x):
        return self.linear(x)

model = MyModel()

output = model(x)
\end{dhkcode}

여기서 \zsh{x}는 적절한 입력 Tensor가 미리 정의되어 있다고 가정한다.

이번 학습을 마친 뒤에는 다음 과정을 스스로 설명할 수 있어야 한다.

\begin{dhknote}
\begin{dhksmry}[centered]
클래스 정의
\end{dhksmry}
\end{dhknote}

\begin{dhkcode}
class MyModel(nn.Module):
\end{dhkcode}

Python이 클래스 객체를 생성하고 \zsh{MyModel}이라는 이름에 바인딩하는 과정

\begin{dhknote}
\begin{dhksmry}[centered]
인스턴스 생성
\end{dhksmry}
\end{dhknote}

\begin{dhkcode}
model = MyModel()
\end{dhkcode}

클래스 객체가 호출되고 인스턴스가 생성되는 과정

\begin{dhknote}
\begin{dhksmry}[centered]
인스턴스 초기화
\end{dhksmry}
\end{dhknote}

\begin{dhkcode}
def __init__(self):
\end{dhkcode}

생성된 인스턴스가 초기화되는 과정

\begin{dhknote}
\begin{dhksmry}[centered]
상속된 기능 활용
\end{dhksmry}
\end{dhknote}

\begin{dhkcode}
super().__init__()
\end{dhkcode}

상속 관계에서 \zsh{super()}가 다음 메서드를 탐색하고 호출하는 과정

\begin{dhknote}
\begin{dhksmry}[centered]
인스턴스 속성 설정
\end{dhksmry}
\end{dhknote}

\begin{dhkcode}
self.linear = nn.Linear(10, 5)
\end{dhkcode}

생성한 객체가 인스턴스에 연결되고 Pytorch에 하위 모듈로 등록되는 과정

\begin{dhknote}
\begin{dhksmry}[centered]
Callable 객체 호출
\end{dhksmry}
\end{dhknote}

\begin{dhkcode}
output = model(x)
\end{dhkcode}

인스턴스 호출이 \zsh{nn.Module}의 호출 메커니즘을 거쳐 \zsh{forward()}를 실행하는 과정

이 과정을 이해하면 PyTorch 모델을 단순히 아키텍처를 구현한 코드로 보는 것이 아니라, Python 객체 모델을 활용하여 구성된 실행가능한 시스템으로 이해할 수 있게 된다.
